\section{定义与假设}

\begin{definition}[计划]
对于转变 $s_0 \to s_1$ 而言，\textbf{计划}是在行动发生之前，对至少一个未来决策窗口、一个进入动作以及一种与目标状态相关的环境布置所作的预先规定。
\end{definition}

\begin{definition}[已解决的进入规格]
令 $U(s_1)$ 表示第~\ref{ch:formal-language} 章中的目标含糊性。若一个计划为进入 $s_1$ 固定了一个可见的第一步动作，从而降低了仍未解决的规格负担，则称这个计划提供了\textbf{已解决的进入规格}。
\end{definition}

\begin{definition}[执行选择复杂度]
对于某个未来决策窗口，令 $D$ 表示在行动发生时仍然需要当场决定的、与行动相关的微决策数量。例子包括：选择地点、选定第一个任务、找到材料，或决定何时结束。
\end{definition}

\begin{remark}
变量 $D$ 被引入为一个定性的记账装置，而不是某个测得的认知常数。下文唯一需要的性质是单调性：如果当场仍有更多决策悬而未决，那么执行负担就更难，而不是更容易。
\end{remark}

\begin{axiom}[单调计划假设]
在保持当前状态不变的情况下，激活障碍 $\DV(s_0 \to s_1)$ 关于目标含糊性弱增，关于相关准备弱减。
\end{axiom}

\begin{axiom}[决策负载假设]
在其他项保持不变时，当场执行负担关于 $D$ 弱增。
\end{axiom}

\begin{discussion}
这些假设并没有证明某一个具体计划者在所有情境中都一定优于某一个具体的自发行动者。它们只是把本书已经在使用的业务逻辑形式化了：未来留下的问题越少，进入就越容易；相关设置提前解决得越多，行动当下的负担就越低。
\end{discussion}